\uuid{JWss}
\titre{Limite en l'infini}
\auteur{bodin ; hanani ; mzari}
\organisation{exo7}
\licence{CC BY-NC-SA 4.0}
\datecreate{2020-01-20}
\niveau{L1}
\module{Analyse}
\chapitre{Continuité, limite et étude de fonctions réelles}
\sousChapitre{Limite de fonctions}
\difficulte{3}
\theme{Limites des fonctions réelles, limite en l'infini}
\exo7qcmid{lille-1.limites.028}

\contenu{
\begin{question}
\qclassification{123.03}
\qtags{motcle=limite en l'infini}
Soit  $f(x)=\sqrt{x^3+2x^2+3}-ax\sqrt{x+b}, \, a,b \in \Rr$. $f$ admet une limite  finie en $+\infty$  si et seulement si  :
\begin{answers}

    \bad{ $ a>0 $ et $b>0$}

    \bad{ $a=1$ et $b>0$ }

    \good{$a=1$ et $b=2$}

    \bad{ $a=1$ et $b=0$}

\end{answers}
\begin{explanations}
Si $a\le 0$, $\lim_{x\to +\infty} f(x)=+\infty$. On suppose donc que $a>0$ et on multiplie $f$ par son expression conjuguée. on obtient :
$f(x)= \frac{(1-a^2)x^3+(2-a^2b)x^2+3}{\sqrt{x^3+2x^2+3}+ax\sqrt{x+b}}$. On déduit que $f$ admet une limite finie en $+\infty$ si et seulement si $a=1$ et $b=2$.
\end{explanations}
\end{question}
}
